
% Phase 2 instructions (appendix).
\section{Phase 2: Design}\label{sec:phase2}
\SetDocHeader{Phase 2: Design}

\subsection{Objective}
The goal of this phase is architectural and interface planning. The team must define the data flow,
security protocols, and user interface layouts. Successful completion of this phase ensures that
the Backend and Frontend teams can work in parallel without integration conflicts.

\subsection{Checkpoints}
\begin{itemize}
  \item[$\square$] \textbf{System Architecture:} Document the data flow (Angular $\rightarrow$ Spring Boot $\rightarrow$ SQLite).
  \item[$\square$] \textbf{Security Planning:} Familiarize yourself with JJWT.
  \item[$\checkmark$] \textbf{UI Wireframing:} Map out the interface for authentication and task management (Section~\ref{sec:wireframes}).
  \item[$\checkmark$] \textbf{Component Design:} Loosely define the Angular component hierarchy and service structure (Section~\ref{sec:task-backlog}, 12 June 2026).
\end{itemize}

\subsection{Technical Reference: Design Standards}

\subsubsection{API Design Principles}
All endpoints must follow RESTful conventions. Use the following templates for your design documentation:


\newpage


\subsubsection{Resource Creation Flow (Template)}
Use this logic to design your \texttt{POST} endpoints to ensure consistent error handling
(e.g., handling duplicate entries).

\begin{figure}[htbp]
  \centering
  \resizebox{\linewidth}{!}{%
  \begin{sequencediagram}
    \newthread{requester}{Client/Requester}
    \newinst[1.3]{api}{API/Interface Layer}
    \newinst[1.3]{logic}{Business Logic Layer}
    \newinst[1.3]{da}{Data Access Layer}
    \newinst[1.3]{store}{Data Storage}
    \begin{sdblock}[gray!10]{SCENARIO A}{[Resource is Unique]}
      \begin{call}{requester}{[Create Request]}{api}{[2xx Success Response]}
        \begin{call}{api}{[Process Creation]}{logic}{[Success Data/Object]}
          \begin{call}[2]{logic}{[Check Existence]}{da}{}
            \begin{call}{da}{[Query]}{store}{[Result (Not Found)]}
            \end{call}
          \end{call}
          \begin{call}[2]{logic}{[Persist New Resource]}{da}{[Resource Object]}
            \begin{call}{da}{[Insert/Write]}{store}{[New Resource Object]}
            \end{call}
          \end{call}
        \end{call}
      \end{call}
    \end{sdblock}
    \begin{sdblock}[gray!10]{SCENARIO B}{[Conflict/Already Exists]}
      \begin{call}{requester}{[Create Request]}{api}{[4xx Error Response]}
        \begin{call}{api}{[Process Creation]}{logic}{[Conflict/Error Signal]}
          \begin{call}[2]{logic}{[Check Existence]}{da}{}
            \begin{call}{da}{[Query]}{store}{[Result (Found)]}
            \end{call}
          \end{call}
        \end{call}
      \end{call}
    \end{sdblock}
  \end{sequencediagram}%
  }
\end{figure}


\newpage



\subsubsection{Login/Auth Flow (Template)}
Use this logic to design your \texttt{/auth} endpoints to ensure secure credential validation.

\begin{figure}[htbp]
  \centering
  \resizebox{\linewidth}{!}{%
  \begin{sequencediagram}
    \newthread{requester}{Client/Requester}
    \newinst[1.1]{api}{API/Interface Layer}
    \newinst[1.1]{logic}{Business Logic Layer}
    \newinst[1.1]{da}{Data Access Layer}
    \newinst[1.1]{util}{Utility/Service}
    \newinst[1.1]{store}{Data Storage}
    \begin{sdblock}[gray!10]{SCENARIO A}{[Success]}
      \begin{call}{requester}{[Login Request]}{api}{[2xx Success Response]}
        \begin{call}{api}{[Execute Process]}{logic}{[Return Success/Token]}
          \begin{call}[2]{logic}{[Request User Data]}{da}{}
            \begin{call}{da}{[Query/Fetch]}{store}{[Return User Object]}
            \end{call}
          \end{call}
          \begin{call}[1]{logic}{[Validate Credentials/Generate Token]}{util}{[Return Result]}
          \end{call}
        \end{call}
      \end{call}
    \end{sdblock}
    \begin{sdblock}[gray!10]{SCENARIO B}{[Failure - Invalid Credentials]}
      \begin{call}{requester}{[Login Request]}{api}{[4xx Error Response]}
        \begin{call}{api}{[Execute Process]}{logic}{[Return Error/Fail]}
          \begin{call}[2]{logic}{[Request User Data]}{da}{}
            \begin{call}{da}{[Query/Fetch]}{store}{[Return Null/Empty]}
            \end{call}
          \end{call}
        \end{call}
      \end{call}
    \end{sdblock}
  \end{sequencediagram}%
  }
\end{figure}



\newpage


\subsection{Deliverables for Phase 2}
By the end of this phase, your team should have:
\begin{enumerate}
  \item \textbf{API Specification:} A completed document defining all endpoints, request bodies, and response codes.
  \item \textbf{Wireframes:} Visual mockups of the Login, Register, and Dashboard screens.
  \item \textbf{Database Schema:} A final relational model (ERD) ready for implementation.
\end{enumerate}

\subsection{Moving to Phase 3}
Once the API contracts and UI wireframes are finalized and agreed upon by all team members,
proceed to \textbf{Phase 3: Backend Development}.

\newpage
\subsection{System Architecture}\label{sec:architecture}
Document the data flow (Angular $\rightarrow$ Spring Boot $\rightarrow$ SQLite).
See deliverable checklist in Section~\ref{sec:phase2}.

\newpage
\subsection{Security Planning}\label{sec:security-planning}
Familiarize yourself with JJWT.
See deliverable checklist in Section~\ref{sec:phase2}.

\newpage
\subsection{UI Wireframing}\label{sec:wireframing}
Map out the interface for authentication and task management.
See deliverable checklist in Section~\ref{sec:phase2}.

\newpage
\subsection{Component Design}\label{sec:component-design}
Loosely define the Angular component hierarchy and service structure.
See deliverable checklist in Section~\ref{sec:phase2}.

\newpage
\subsection{API Specification}\label{sec:api-spec}
Complete the document defining all endpoints, request bodies, and response codes.
See deliverable checklist in Section~\ref{sec:phase2}.

\newpage
\subsection{Database Schema}\label{sec:database-schema}
Finalize the relational model (ERD) ready for implementation.
See deliverable checklist in Section~\ref{sec:phase2}.

\newpage
